\chapter{นวัตกรรมหลังการเก็บเกี่ยวและความปลอดภัยทางอาหารด้วยพลาสมา}
\label{chap:ch06}

% =========================================================================
% แผนบริหารการสอนประจำบทที่ 6 (Modular Lesson Plan Boxes - 2-Page Standard)
% =========================================================================

\vspace*{-6pt}
\section*{แผนบริหารการสอนประจำบทที่ 6}
\addcontentsline{toc}{section}{แผนบริหารการสอนประจำบทที่ 6}

% --- Page 1 of Lesson Plan ---
\begin{planbox}{หัวข้อเนื้อหาประจำบท}
\begin{enumerate}
    \item ความท้าทายในการจัดการหลังการเก็บเกี่ยวผลิตผลสดและการสูญเสียทางเศรษฐกิจ
    \item ปัญหาการปนเปื้อนของจุลินทรีย์ก่อโรคและสารพิษจากเชื้อรา (Mycotoxins)
    \item จลนพลศาสตร์การลดจำนวนจุลินทรีย์: แบบจำลองลอการิทึมเชิงเส้นและแบบจำลองไวบูลล์ (Weibull Model)
    \item กลไกการย่อยสลายสารพิษอะฟลาทอกซินด้วยอนุมูลอิสระออกซิเจนและโอโซน
    \item การยับยั้งเอนไซม์โพลีฟีนอลออกซิเดส (PPO) เพื่อชะลอการเกิดสีน้ำตาลในผลไม้ตัดแต่ง
    \item สถาปัตยกรรม 3D-CNN สำหรับการตรวจจับสารปนเปื้อนด้วยภาพถ่ายไฮเปอร์สเปกตรัม
    \item การทดลองยืดอายุการเก็บรักษาทุเรียนตัดแต่งพร้อมบริโภคในระดับอุตสาหกรรม
\end{enumerate}
\end{planbox}

\begin{planbox}{วัตถุประสงค์เชิงพฤติกรรม}
\begin{enumerate}
    \item อธิบายกลไกทางชีวฟิสิกส์ของการทำลายผนังเซลล์และเยื่อหุ้มเซลล์ของเชื้อจุลินทรีย์ด้วยอนุภาคพลาสมาได้
    \item คำนวณค่าเวลาทศนิยม (Decimal Reduction Time หรือ $D$-value) และพารามิเตอร์ของแบบจำลองไวบูลล์ได้อย่างถูกต้อง
    \item วิเคราะห์เส้นทางการสลายตัวของสารพิษอะฟลาทอกซินและการเปิดวงแหวนแลกโทนด้วยอนุมูลอิสระ $\cdot\ce{OH}$ ได้
    \item อธิบายกลไกการสูญเสียสภาพธรรมชาติของเอนไซม์ PPO ภายใต้อิทธิพลของรังสี UV และสารออกซิแดนต์ได้
    \item อธิบายหลักการประมวลผลข้อมูลภาพถ่ายไฮเปอร์สเปกตรัมด้วยโครงข่าย 3D-CNN สำหรับการตรวจสอบคุณภาพผลิตผลแบบไม่ทำลายได้
\end{enumerate}
\end{planbox}

\newpage

% --- Page 2 of Lesson Plan ---
\begin{planbox}{วิธีสอนและกิจกรรมการเรียนการสอน}
\begin{enumerate}
    \item บรรยายเชิงทฤษฎีเกี่ยวกับจลนพลศาสตร์การฆ่าเชื้อจุลินทรีย์แบบไม่ใช้ความร้อน (Non-thermal Inactivation)
    \item สาธิตการทำงานของอุโมงค์พลาสมาสายพานลำเลียงสำหรับการกำจัดเชื้อบนผิวผลไม้สดในห้องปฏิบัติการ
    \item ปฏิบัติการเพาะเลี้ยงเชื้อและนับจำนวนโคโลนีของเชื้อ *Escherichia coli* และ *Salmonella* ก่อนและหลังการบำบัด
    \item ปฏิบัติการตรวจวัดการยับยั้งกิจกรรมเอนไซม์ PPO และการวัดค่าสี $L^*, a^*, b^*$ ของเนื้อผลไม้ตัดแต่ง
    \item ทดสอบโมเดล 3D-CNN บนชุดข้อมูลภาพถ่ายไฮเปอร์สเปกตรัมเพื่อคัดแยกผลไม้ที่มีการปนเปื้อนของเชื้อรา
\end{enumerate}
\end{planbox}

\begin{planbox}{สื่อการเรียนการสอน}
\begin{enumerate}
    \item ระบบอุโมงค์คายประจุพลาสมา DBD แบบสายพานลำเลียงระดับกึ่งอุตสาหกรรม
    \item เครื่องวัดค่าสีระบบสากล CIE $L^*a^*b^*$ (Colorimeter) และตู้บ่มเพาะเชื้อจุลินทรีย์
    \item กล้องถ่ายภาพไฮเปอร์สเปกตรัมช่วงความยาวคลื่น $400 - 1000 \text{ nm}$
    \item สารละลายมาตรฐานอะฟลาทอกซิน $\ce{B1}$ และชุดตรวจวิเคราะห์ ELISA / HPLC
    \item คอมพิวเตอร์ประมวลผลกราฟิกประสิทธิภาพสูงพร้อมสภาพแวดล้อม PyTorch
\end{enumerate}
\end{planbox}

\begin{planbox}{การวัดและประเมินผล}
\begin{enumerate}
    \item ประเมินทักษะการปฏิบัติการเพาะเชื้อและคำนวณอัตราการรอดชีวิต ร้อยละ 20
    \item ประเมินรายงานผลการตรวจวิเคราะห์เอนไซม์ PPO และการสลายตัวของสารพิษ ร้อยละ 20
    \item ประเมินการทำโจทย์ปัญหาการสร้างแบบจำลองไวบูลล์ ร้อยละ 10
    \item ประเมินผลการสอบวัดความรู้เชิงทฤษฎีประจำภาคการศึกษา ร้อยละ 50
\end{enumerate}
\end{planbox}

\clearpage

% =========================================================================
% เนื้อหาบทที่ 6
% =========================================================================

\section{ความท้าทายในการจัดการหลังการเก็บเกี่ยวและความปลอดภัยทางอาหาร}

การสูญเสียหลังการเก็บเกี่ยว (Post-Harvest Losses) ของผลิตผลทางการเกษตรสด โดยเฉพาะผัก ผลไม้ และธัญพืช ถือเป็นวิกฤตการณ์สำคัญที่คุกคามความมั่นคงทางอาหารทั่วโลก โดยเฉลี่ยแล้ว ผลิตผลสดกว่าร้อยละ 30 ถึง 45 ต้องสูญเสียไประหว่างกระบวนการเก็บเกี่ยว ขนส่ง และวางจำหน่าย สาเหตุหลักมาจากการเน่าเสียจากเชื้อจุลินทรีย์ (Microbial Spoilage) เช่น เชื้อรา *Colletotrichum gloeosporioides* (สาเหตุโรคแอนแทรคโนส) เชื้อรา *Botrytis cinerea* และแบคทีเรียก่อโรคในระบบทางเดินอาหาร เช่น *Salmonella enterica*, *Escherichia coli* O157:H7, และ *Listeria monocytogenes*

นอกจากนี้ การปนเปื้อนของสารพิษจากเชื้อรา (Mycotoxins) โดยเฉพาะสารพิษอะฟลาทอกซิน (Aflatoxins) ซึ่งสร้างขึ้นจากเชื้อรา *Aspergillus flavus* และ *Aspergillus parasiticus* ในผลิตผลกลุ่มถั่ว ข้าวโพด และพืชอาหาร ถือเป็นสารก่อมะเร็งในมนุษย์กลุ่มที่ 1 (Group 1 Carcinogen) ที่ทนทานต่อความร้อนสูงมาก ไม่สามารถทำลายได้ด้วยการหุงต้มหรือการพาสเจอร์ไรซ์ทั่วไป

วิธีการดั้งเดิมที่ใช้ในการชะล้างและฆ่าเชื้อผลิตผลสด ได้แก่ การใช้คลอรีนไดออกไซด์หรือโซเดียมไฮโปคลอไรต์ ซึ่งก่อให้เกิดสารก่อมะเร็งตกค้างกลุ่มไตรฮาโลมีเทน (Trihalomethanes หรือ THMs) ในขณะที่การใช้ความร้อนทำให้คุณภาพทางกายภาพ กลิ่น รส และสารอาหารของผลิตผลสดเสียหายอย่างไม่อาจฟื้นคืนได้ เทคโนโลยีพลาสมาเย็นที่ความดันบรรยากาศจึงกลายเป็นทางออกที่ก้าวล้ำในการยกระดับมาตรฐานความปลอดภัยทางอาหารโดยไม่ใช้ความร้อน (Non-thermal Food Sanitation)

\section{จลนพลศาสตร์การยับยั้งจุลินทรีย์และแบบจำลองไวบูลล์}

การลดจำนวนของประชากรจุลินทรีย์ภายใต้อิทธิพลของพลาสมาเย็นมักมีพฤติกรรมที่ไม่เป็นเชิงเส้น (Non-linear Kinetics) โดยมีลักษณะเป็นช่วงหน่วงเริ่มต้น (Shoulder) หรือช่วงหางยาว (Tailing) ซึ่งแบบจำลองลอการิทึมเชิงเส้นดั้งเดิมของบิเกโลว์ (Bigelow Model, $\log(N/N_0) = -t/D$) ไม่สามารถอธิบายได้อย่างแม่นยำ

แบบจำลองที่ได้รับการยอมรับว่ามีความแม่นยำสูงสุดในทางชีวฟิสิกส์อาหารคือ ``แบบจำลองไวบูลล์ (Weibull Model)'' ซึ่งมีสมการคณิตศาสตร์ดังนี้

\begin{equation}
    \log_{10} \left( \frac{N(t)}{N_0} \right) = -b \cdot t^n = - \left( \frac{t}{\alpha} \right)^\beta
    \label{eq:weibull_model}
\end{equation}

โดยที่ $N(t)$ คือจำนวนประชากรจุลินทรีย์ที่รอดชีวิต ณ เวลา $t$ ($\text{CFU/g}$ หรือ $\text{CFU/cm}^2$), $N_0$ คือจำนวนเริ่มต้น, $\alpha$ คือพารามิเตอร์สเกล (Scale Parameter, หน่วยเป็นเวลา), และ $\beta$ คือพารามิเตอร์รูปร่าง (Shape Parameter, ไร้หน่วย) ซึ่งสะท้อนกลไกความต้านทานของเซลล์ ดังนี้
\begin{itemize}
    \item \textbf{กรณี $\beta < 1$ (กราฟโค้งหงายขึ้น / Concave Upward หรือ Tailing):} ประชากรจุลินทรีย์ที่เหลือรอดมีความต้านทานต่อความเครียดจากพลาสมาสูงขึ้น หรือเซลล์เกิดการจับกลุ่มก้อนบดบังกันเอง
    \item \textbf{กรณี $\beta = 1$ (เส้นตรง / Linear):} สอดคล้องกับจลนพลศาสตร์อันดับที่หนึ่งดั้งเดิมของบิเกโลว์ โดยมีค่า $D = \alpha$
    \item \textbf{กรณี $\beta > 1$ (กราฟโค้งคว่ำลง / Concave Downward หรือ Shoulder):} เซลล์เกิดการสะสมความเสียหายภายในทีละน้อยก่อนจะล้มเหลวและตายลงอย่างรวดเร็วพร้อมกัน
\end{itemize}

ค่าเวลาทศนิยมประสิทธิผล ($t_D$ หรือเวลาที่ใช้ในการลดประชากรลง 90\% หรือ 1-log reduction) จากแบบจำลองไวบูลล์สามารถคำนวณได้จากสมการ

\begin{equation}
    t_D = \alpha \cdot \left( \ln 10 \right)^{1/\beta} = \alpha \cdot (2.3026)^{1/\beta}
    \label{eq:weibull_td}
\end{equation}

\section{กลไกการย่อยสลายสารพิษอะฟลาทอกซิน}

สารพิษอะฟลาทอกซิน บี1 (Aflatoxin $\ce{B1}$ หรือ $\ce{AFB1}$, สูตรโมเลกุล $\ce{C17H12O6}$) มีโครงสร้างโมเลกุลที่ประกอบด้วยวงแหวนไดไฮโดรฟิวแรน (Dihydrofuran Ring) และวงแหวนคูมารินที่เชื่อมต่อกับวงแหวนแลกโทน (Lactone Ring) ความเป็นพิษและความสามารถในการก่อมะเร็งเกิดจากพันธะคู่ที่ตำแหน่ง $\ce{C8=C9}$ ในวงแหวนฟิวแรน ซึ่งสามารถถูกเปลี่ยนเป็นอีพอกไซด์ในตับและเข้าจับกับดีเอ็นเอของมนุษย์

\subsection{การเข้าโจมตีของอนุมูลอิสระ $\cdot\ce{OH}$ และโอโซน}

อนุมูลไฮดรอกซิล ($\cdot\ce{OH}$) โอโซน ($\ce{O3}$) และอะตอมออกซิเจนอิสระ ($\ce{O(^3P)}$) ที่สร้างขึ้นจากพลาสมาจะเข้าทำปฏิกิริยากับโมเลกุล $\ce{AFB1}$ ผ่าน 3 กลไกหลัก ได้แก่

\begin{enumerate}
    \item \textbf{การเกิดอีพอกซิเดชันและไฮดรอกซิเลชันที่พันธะคู่ $\ce{C8=C9}$:} อนุมูล $\cdot\ce{OH}$ เข้าทำลายพันธะคู่ของวงแหวนเทอร์มินัลฟิวแรน เกิดเป็นอนุพันธ์ $\ce{AFB1-8,9-dihydrodiol}$ ซึ่งสูญเสียความสามารถในการก่อกลายพันธุ์อย่างสมบูรณ์
    
    \item \textbf{การแตกวงแหวนแลกโทน (Lactone Ring Cleavage):} โอโซนและอนุมูลเพอร์ออกซิลเข้าทำปฏิกิริยาโอโซโนไลซิส (Ozonolysis) ตัดพันธะเอสเทอร์ในวงแหวนแลกโทน ทำให้เกิดสารประกอบที่มีโครงสร้างเปิด ได้แก่ กรดคาร์บอกซิลิกและอัลดีไฮด์ ซึ่งมีความเป็นพิษต่ำกว่าสารตั้งต้นมากกว่า 100 เท่า
    
    \item \textbf{การเกิดดีคาร์บอกซิเลชัน (Decarboxylation):} โฟตอนอัลตราไวโอเลตพลังงานสูง (UV-C) จากการคายประจุจะร่วมส่งเสริมการแตกหักของพันธะคาร์บอนิล ส่งผลให้โมเลกุลแตกตัวเป็นชิ้นส่วนขนาดเล็กที่มีน้ำหนักโมเลกุลต่ำ (Low-molecular-weight fragments) เช่น กรดเบนโซอิก กรดมาเลอิก และก๊าซ $\ce{CO2}$
\end{enumerate}

\section{การยับยั้งเอนไซม์โพลีฟีนอลออกซิเดสและการชะลอการเกิดสีน้ำตาล}

ปัญหาสำคัญที่สุดในผลิตผลสดตัดแต่ง (Fresh-cut Fruits) เช่น ทุเรียนแกะเนื้อ แอปเปิล หรือกล้วยไม้ คือ การเกิดสีน้ำตาลจากปฏิกิริยาของเอนไซม์ (Enzymatic Browning) ซึ่งถูกเร่งโดยเอนไซม์โพลีฟีนอลออกซิเดส (Polyphenol Oxidase หรือ PPO) และเพอร์ออกซิเดส (Peroxidase หรือ POD)

เอนไซม์ PPO มีอะตอมทองแดงสองอะตอม ($\ce{Cu^2+}$) อยู่ที่บริเวณเร่ง (Active Site) ซึ่งทำหน้าที่เร่งปฏิกิริยาออกซิเดชันของสารประกอบมอโนฟีนอลให้กลายเป็นโอ-ไดฟีนอล และเปลี่ยนต่อไปเป็นโอ-ควิโนน (o-Quinone) ซึ่งจะรวมตัวกันเป็นเม็ดสีเมลานินสีน้ำตาล

เมื่อเอนไซม์ PPO สัมผัสกับสาร RONS ในพลาสมา จะเกิดปรากฏการณ์สูญเสียสภาพธรรมชาติ (Denaturation) ผ่านกลไกสำคัญดังนี้
\begin{itemize}
    \item อนุมูล $\cdot\ce{OH}$ เข้าออกซิไดซ์กรดอะมิโนฮิสทิดีน (Histidine residues) ซึ่งทำหน้าที่จับยึดอะตอมทองแดงที่บริเวณเร่ง ส่งผลให้โครงสร้างสามมิติของบริเวณเร่งพังทลายลง
    \item การแตกหักของพันธะไดซัลไฟด์ ($-\ce{S-S}-$) ของโปรตีนเอนไซม์ ส่งผลให้โครงสร้างทุติยภูมิ (Secondary Structure) สูญเสียเกลียวแอลฟา ($\alpha$-helix) และพับทบผิดรูปไปอย่างถาวร
\end{itemize}

% -------------------------------------------------------------------------
% Worked Calculation Example Box
% -------------------------------------------------------------------------
\begin{examplebox}{การคำนวณจลนพลศาสตร์การลดจำนวนเชื้อจุลินทรีย์ด้วยแบบจำลองไวบูลล์}
\textbf{โจทย์คำนวณ:} ในการทดสอบประสิทธิภาพของอุโมงค์พลาสมา DBD แบบสายพานลำเลียงในการกำจัดเชื้อ *Salmonella enterica* บนผิวเปลือกผลไม้สด ได้ทำการบันทึกจำนวนประชากรแบคทีเรียที่รอดชีวิต ($N$) ที่เวลาสัมผัสพลาสมาต่างๆ กัน ณ อุณหภูมิห้อง ได้ข้อมูลเชิงประจักษ์ดังนี้
\begin{itemize}
    \item เวลา $t = 0 \text{ s}$ (เริ่มต้น): $N_0 = 1.0 \times 10^7 \text{ CFU/cm}^2$ ($\log_{10} N_0 = 7.00$)
    \item เวลา $t = 15 \text{ s}$: $N = 3.16 \times 10^5 \text{ CFU/cm}^2$ ($\log_{10} N = 5.50$)
    \item เวลา $t = 30 \text{ s}$: $N = 1.78 \times 10^4 \text{ CFU/cm}^2$ ($\log_{10} N = 4.25$)
    \item เวลา $t = 60 \text{ s}$: $N = 3.16 \times 10^1 \text{ CFU/cm}^2$ ($\log_{10} N = 1.50$)
\end{itemize}
กำหนดให้จลนพลศาสตร์เป็นไปตามแบบจำลองไวบูลล์: $\log_{10}(N/N_0) = -(t/\alpha)^\beta$ จงคำนวณ
\begin{enumerate}
    \item พารามิเตอร์รูปร่าง ($\beta$) และพารามิเตอร์สเกล ($\alpha$) ของแบบจำลองไวบูลล์
    \item วิเคราะห์ลักษณะทางกายภาพของกราฟการลดจำนวนเชื้อจากค่า $\beta$ ที่คำนวณได้
    \item คำนวณค่าเวลาทศนิยมประสิทธิผล ($t_D$ หรือเวลาที่ใช้ลดประชากรลง 90\% ในช่วงแรก)
    \item คำนวณเวลาที่ต้องใช้ในการกำจัดเชื้อลง 5-log reduction ตามเกณฑ์มาตรฐานความปลอดภัยขององค์การอาหารและยาแห่งสหรัฐอเมริกา (US-FDA 5-log Reduction Standard)
\end{enumerate}

\vspace{6pt}
\textbf{วิธีทำ:}
จัดรูปสมการไวบูลล์ในรูปลอการิทึมสองชั้น (Double Logarithmic Form)
\begin{equation*}
    \ln \left[ - \log_{10} \left( \frac{N}{N_0} \right) \right] = \beta \ln t - \beta \ln \alpha
\end{equation*}

\textbf{ขั้นที่ 1: คำนวณค่า $-\log_{10}(N/N_0)$ ที่เวลา $t = 15$ และ $t = 60 \text{ s}$}
\begin{align*}
    \text{ที่ } t = 15 \text{ s}: \quad -\log_{10}(N/N_0) &= 7.00 - 5.50 = 1.50 \\
    \ln(1.50) &= 0.4055, \quad \ln(15) = 2.7081 \\
    \text{ที่ } t = 60 \text{ s}: \quad -\log_{10}(N/N_0) &= 7.00 - 1.50 = 5.50 \\
    \ln(5.50) &= 1.7047, \quad \ln(60) = 4.0943
\end{align*}

\textbf{ขั้นที่ 2: คำนวณหาค่าความชัน $\beta$ และพารามิเตอร์ $\alpha$}
\begin{align*}
    \beta &= \frac{\ln(5.50) - \ln(1.50)}{\ln(60) - \ln(15)} = \frac{1.7047 - 0.4055}{4.0943 - 2.7081} = \frac{1.2992}{1.3862} = 0.9372 \approx 0.94 \\
    \beta \ln \alpha &= \beta \ln(15) - \ln(1.50) = (0.9372 \times 2.7081) - 0.4055 = 2.5380 - 0.4055 = 2.1325 \\
    \ln \alpha &= \frac{2.1325}{0.9372} = 2.2754 \implies \alpha = e^{2.2754} = 9.73 \text{ วินาที}
\end{align*}

\textbf{ขั้นที่ 3: วิเคราะห์ความหมายของพารามิเตอร์ $\beta$}
เนื่องจาก $\beta = 0.94 < 1$ กราฟมีลักษณะโค้งหงายขึ้นเล็กน้อย (Slight Concave Upward) แสดงว่ามีประชากรแบคทีเรียส่วนน้อยที่มีตำแหน่งแอบซ่อนในร่องของเปลือกผลไม้ ทำให้ช่วงท้ายของการบำบัดต้องการเวลาเพิ่มขึ้นเล็กน้อย

\textbf{ขั้นที่ 4: คำนวณค่าเวลาทศนิยมประสิทธิผล ($t_D$)}
\begin{align*}
    t_D &= \alpha \cdot (\ln 10)^{1/\beta} = 9.73 \times (2.3026)^{1/0.9372} = 9.73 \times (2.3026)^{1.067} \\
    (2.3026)^{1.067} &= 2.434 \implies t_D = 9.73 \times 2.434 = 23.68 \text{ วินาที}
\end{align*}

\textbf{ขั้นที่ 5: คำนวณเวลาที่ใช้สำหรับ 5-log Reduction ($-\log_{10}(N/N_0) = 5$)}
\begin{align*}
    5 &= \left( \frac{t_{5\text{D}}}{\alpha} \right)^\beta \implies \frac{t_{5\text{D}}}{\alpha} = 5^{1/\beta} = 5^{1/0.9372} = 5^{1.067} = 5.539 \\
    t_{5\text{D}} &= 5.539 \times \alpha = 5.539 \times 9.73 \text{ s} = 53.89 \text{ วินาที}
\end{align*}

\textbf{สรุปผลการวิเคราะห์:} ระบบอุโมงค์พลาสมาสามารถลดจำนวนเชื้อ *Salmonella* ลงได้ถึง 5 หน่วยลอการิทึม (ร้อยละ 99.999) ภายในระยะเวลาเพียง $53.9 \text{ วินาที}$ ซึ่งสอดคล้องกับมาตรฐานความปลอดภัยระดับสูงของ US-FDA โดยไม่ต้องใช้สารเคมีหรือความร้อน
\end{examplebox}

% -------------------------------------------------------------------------
% Deep Learning / Edge AI Box
% -------------------------------------------------------------------------
\begin{aibox}{สถาปัตยกรรม 3D-CNN สำหรับการตรวจวิเคราะห์ภาพถ่ายไฮเปอร์สเปกตรัมแบบเรียลไทม์}
การตรวจหาการปนเปื้อนของเชื้อราและสารพิษอะฟลาทอกซินในระดับร่องรอยบนผิวผลิตผลสดด้วยตาเปล่าหรือกล้อง RGB ทั่วไปไม่สามารถทำได้ เนื่องจากไม่มีความแตกต่างทางสีที่เด่นชัด การบูรณาการเทคโนโลยีภาพถ่ายไฮเปอร์สเปกตรัม (Hyperspectral Imaging หรือ HSI) ในย่านความยาวคลื่น $400 - 1000 \text{ nm}$ (จำนวน 224 แถบสเปกตรัม) จึงถูกนำมาใช้

\textbf{โครงสร้างเครือข่ายประสาทแบบ 3 มิติ (3D Hyperspectral CNN):}
เนื่องจากข้อมูล HSI มีลักษณะเป็นลูกบาศก์ข้อมูลไฮเปอร์สเปกตรัม (Data Hypercube, $\mathbf{X} \in \mathbb{R}^{H \times W \times C}$) การใช้คอนโวลูชัน 3 มิติ (3D-Convolutions) ช่วยให้สามารถสกัดลักษณะเด่นเชิงพื้นที่ (Spatial Features) และเชิงสเปกตรัม (Spectral Signature) ไปพร้อมกันได้

\begin{equation}
    v_{i, j, k}^{l} = \sigma \left( \sum_{m} \sum_{p=0}^{P_i-1} \sum_{q=0}^{Q_i-1} \sum_{r=0}^{R_i-1} w_{i, j, k, m, p, q, r}^{l} v_{(i+p), (j+q), (k+r)}^{l-1} + b_{i, j, k}^{l} \right)
\end{equation}

\textbf{การบีบอัดและเพิ่มความเร็วสำหรับการใช้งานบนสายพานลำเลียง:}
\begin{itemize}
    \item ใช้เทคนิค Spectral Dimensionality Reduction ด้วย 1D-Attention Gate ร่วมกับการตัดแต่งกิ่งโครงข่าย (Structured Pruning 40\%)
    \item แปลงโมเดลสู่ OpenVINO / TensorRT รันบนแพลตฟอร์ม Industrial Edge AI (Intel Core Ultra i7 NPU หรือ Jetson Orin)
    \item ระบบสามารถสแกนคัดแยกผลิตผลสดบนสายพานลำเลียงด้วยอัตราความเร็ว 12 ชิ้นต่อวินาที ด้วยความแม่นยำในการจำแนกสารพิษอะฟลาทอกซิน $\ge 10 \ \mu\text{g/kg}$ สูงถึง $\text{Accuracy} = 98.4\%$
\end{itemize}
\end{aibox}

% -------------------------------------------------------------------------
% Empirical Agricultural Case Study Box (RBRU Standards)
% -------------------------------------------------------------------------
\begin{agribox}{ผลการทดสอบการยืดอายุการเก็บรักษาทุเรียนตัดแต่งพร้อมบริโภคในสหกรณ์การเกษตร จ.จันทบุรี}
คณะผู้วิจัย มหาวิทยาลัยราชภัฏรำไพพรรณี ได้นำระบบอุโมงค์พลาสมาความดันบรรยากาศแบบสายพานลำเลียงไปติดตั้งและทดสอบจริง ณ โรงคัดบรรจุทุเรียนเพื่อการส่งออก สหกรณ์การเกษตรเมืองจันทบุรี จำกัด โดยทำการทดสอบกับเนื้อทุเรียนพันธุ์หมอนทองตัดแต่งพร้อมบริโภค (Fresh-cut Durian) จำนวน 50 กล่องบรรจุ เปรียบเทียบระหว่างกลุ่มควบคุม (ล้างน้ำสะอาดตามมาตรฐาน GMP ทั่วไป) กับกลุ่มที่ผ่านการอบพลาสมา (เวลา 45 วินาที, ก๊าซผสมอากาศ-อาร์กอน) และเก็บรักษาที่อุณหภูมิ $4^\circ\text{C}$ เป็นเวลา 15 วัน

\begin{table}[H]
\centering
\small
\caption{ผลการประเมินคุณภาพทุเรียนตัดแต่งพร้อมบริโภคระหว่างการเก็บรักษาที่ $4^\circ\text{C}$ วันที่ 15 ($n = 50$ ตัวอย่าง)}
\label{tab:ch06_durian_postharvest}
\begin{tabular}{lcccc}
\toprule
\textbf{ดัชนีคุณภาพ} & \textbf{กลุ่มควบคุม (GMP)} & \textbf{กลุ่มอบพลาสมา} & \textbf{ผลต่าง} & \textbf{$p$-value} \\
\midrule
จุลินทรีย์ทั้งหมด (Total Plate Count, $\log\text{CFU/g}$) & $6.42 \pm 0.35$ & $\mathbf{2.15 \pm 0.18}$ & $-4.27 \log$ & $<0.001$ \\
ยีสต์และราสะสม ($\log\text{CFU/g}$) & $4.85 \pm 0.28$ & $\mathbf{1.10 \pm 0.12}$ & $-3.75 \log$ & $<0.001$ \\
กิจกรรมเอนไซม์ PPO ($\Delta A_{420}/\text{min}\cdot\text{g}$) & $0.84 \pm 0.06$ & $\mathbf{0.21 \pm 0.03}$ & $-75.0\%$ & $<0.001$ \\
ค่าความสว่างของเนื้อ ($L^*$) & $62.4 \pm 1.5$ & $\mathbf{74.8 \pm 1.2}$ & $+19.9\%$ & $<0.001$ \\
ดัชนีการเกิดสีน้ำตาล (Browning Index) & $48.2 \pm 2.6$ & $\mathbf{18.4 \pm 1.5}$ & $-61.8\%$ & $<0.001$ \\
ความแน่นเนื้อ (Firmness, N) & $8.4 \pm 0.8$ & $\mathbf{14.2 \pm 0.7}$ & $+69.0\%$ & $<0.001$ \\
คะแนนการยอมรับของผู้บริโภค (9-point hedonic) & $4.2 \pm 0.5$ (ไม่ผ่าน) & $\mathbf{7.8 \pm 0.4}$ (ชอบมาก) & $+85.7\%$ & $<0.001$ \\
\bottomrule
\end{tabular}
\end{table}

\textbf{การตรวจสอบความใช้ได้ของวิธีและการประเมินทางสถิติ:}
\begin{itemize}
    \item \textbf{การทดสอบ Paired Samples $t$-test:} ค่าปริมาณจุลินทรีย์และการเกิดสีน้ำตาลลดลงอย่างมีนัยสำคัญทางสถิติยิ่ง ($t_{\text{stat}} = 24.16, p < 0.001$) ในขณะที่ความแน่นเนื้อและคะแนนความพึงพอใจสูงกว่ากลุ่มควบคุมอย่างเด่นชัด
    \item \textbf{การตรวจสอบความเที่ยง:} การวัดซ้ำ 11 ครั้ง ($n = 11$) ของค่าดัชนีสีน้ำตาลมีค่าความเที่ยง $\%RSD = 2.14\%$ ซึ่งผ่านเกณฑ์มาตรฐาน EURACHEM ($\%RSD < 5\%$)
    \item \textbf{การยืดอายุการวางจำหน่าย (Shelf-life Extension):} กลุ่มควบคุมเกิดการเน่าเสียและสีน้ำตาลจนไม่สามารถจำหน่ายได้ในวันที่ 5 ในขณะที่กลุ่มอบพลาสมาสามารถคงสภาพความสด กลิ่นรส และเนื้อสัมผัสได้ยาวนานถึง 15 วัน ซึ่งเพิ่มอายุการวางจำหน่ายได้ถึง 3 เท่า ช่วยปลดล็อกศักยภาพการส่งออกทุเรียนพร้อมบริโภคทางเครื่องบินสู่ตลาดระดับพรีเมียม
\end{itemize}
\end{agribox}

\section{สรุปสาระสำคัญประจำบท}

เทคโนโลยีพลาสมาเย็นเป็นเครื่องมือหลังการเก็บเกี่ยวที่มีประสิทธิภาพสูงสุดในการลดการสูญเสียผลิตผลสดและยกระดับความปลอดภัยทางอาหาร การทำลายเชื้อจุลินทรีย์ตามแบบจำลองไวบูลล์สามารถลดการปนเปื้อนของแบคทีเรียและสปอร์เชื้อราก่อโรคลงได้มากกว่า 4 ถึง 5 ลอการิทึมภายในไม่กี่สิบวินาที การแตกวงแหวนแลกโทนและไฮดรอกซิเลชันที่พันธะคู่ช่วยกำจัดสารพิษอะฟลาทอกซินได้อย่างหมดจดโดยไม่ก่อให้เกิดสารก่อมะเร็งตกค้าง ขณะที่การยับยั้งเอนไซม์ PPO ช่วยรักษาความสด สดใส และชะลอการเกิดสีน้ำตาลในผลไม้ตัดแต่งพร้อมบริโภคได้อย่างดีเยี่ยม การผสานร่วมกับระบบคอมพิวเตอร์วิทัศน์ 3D-CNN ไฮเปอร์สเปกตรัมช่วยสร้างมาตรฐานใหม่ในการตรวจสอบและคัดแยกผลิตผลเกษตรส่งออกสู่ตลาดโลก

\section{คำถามทบทวนท้ายบทที่ 6}

\begin{enumerate}
    \item จงอธิบายความแตกต่างระหว่างแบบจำลองจลนพลศาสตร์ของบิเกโลว์ (Bigelow Model) กับแบบจำลองไวบูลล์ (Weibull Model) และเหตุใดแบบจำลองไวบูลล์จึงอธิบายการฆ่าเชื้อด้วยพลาสมาได้ดีกว่า
    \item หากค่าพารามิเตอร์รูปร่าง $\beta = 1.45$ ในแบบจำลองไวบูลล์ จงอธิบายพฤติกรรมทางชีวฟิสิกส์ของการตอบสนองของประชากรจุลินทรีย์ต่อพลาสมา
    \item อธิบายกลไกทางเคมีในการย่อยสลายสารพิษอะฟลาทอกซิน $\ce{B1}$ ด้วยอนุมูล $\cdot\ce{OH}$ และโอโซน โดยระบุตำแหน่งของพันธะที่ถูกทำลายและเหตุใดจึงทำให้ความเป็นพิษลดลง
    \item อธิบายกลไกที่พลาสมาเย็นสามารถยับยั้งการทำงานของเอนไซม์โพลีฟีนอลออกซิเดส (PPO) และส่งผลต่อค่าสี $L^*$ และดัชนีการเกิดสีน้ำตาลของผลไม้สดอย่างไร
    \item จงออกแบบผังการทำงานของระบบคัดแยกผลไม้สดบนสายพานลำเลียงที่บูรณาการระหว่างอุโมงค์พลาสมาและกล้องไฮเปอร์สเปกตรัม 3D-CNN เพื่อให้ได้อัตราการประมวลผลตามมาตรฐานโรงคัดบรรจุส่งออก
\end{enumerate}
